\unitopen{37}
\begin{candos}
\item rozpoznawać i tworzyć styl nominalny w oficjalnym piśmie (umowa, regulamin, komunikat);
\item przekształcać zdania z czasownikiem w grupy rzeczownikowe i odwrotnie;
\item używać najczęstszych \de{Funktionsverbgefüge} (\de{eine Entscheidung treffen, zur Verfügung stellen});
\item rozumieć długie zdania urzędowe, dzieląc je na części.
\end{candos}

\section{Słownictwo}
\vocabtable{37}

\section{Czasownik na rzeczownik}
Pisemny niemiecki zawodowy lubi rzeczowniki: zamiast \emph{co ktoś robi} piszesz \emph{co zostaje zrobione}. Nominalizacja dodaje
\de{-ung}, \de{-heit}, \de{-keit}, \de{-ion} albo bierze bezokolicznik z rodzajnikiem (\de{das Lesen}).

\begin{center}\small
\begin{tabular}{@{}p{0.45\linewidth}p{0.47\linewidth}@{}}
\toprule
\textsf{\textbf{Styl czasownikowy}} & \textsf{\textbf{Styl nominalny}}\\
\midrule
\de{Wir prüfen den Antrag.} & \de{Die Prüfung des Antrags erfolgt durch uns.}\\
\de{Wenn man den Fehler behebt, \dots} & \de{Nach der Behebung des Fehlers \dots}\\
\de{Weil das Wetter schlecht war, \dots} & \de{Wegen des schlechten Wetters \dots}\\
\de{Man kann das System nutzen, nachdem man es eingeführt hat.} & \de{Die Nutzung des Systems ist nach seiner Einführung möglich.}\\
\bottomrule
\end{tabular}
\end{center}

\section{\de{Funktionsverbgefüge}}
Czasownik ogólny + rzeczownik zamiast czasownika pełnoznacznego, częste w piśmie urzędowym:

\begin{center}\small
\begin{tabular}{@{}ll@{}}
\toprule
\textsf{\textbf{Funktionsverbgefüge}} & \textsf{\textbf{Prosty odpowiednik}}\\
\midrule
\de{eine Entscheidung treffen} & \de{entscheiden}\\
\de{zur Verfügung stellen} & \de{bereitstellen}\\
\de{in Betrieb nehmen} & \de{starten}\\
\de{Kritik üben} & \de{kritisieren}\\
\de{in Anspruch nehmen} & \de{nutzen}\\
\de{zur Anwendung kommen} & \de{angewendet werden}\\
\bottomrule
\end{tabular}
\end{center}

\begin{zasada}[Kiedy styl nominalny?]
Dokumenty prawne, regulaminy, protokoły i raporty. W mowie i w wiadomościach w zespole preferuj styl czasownikowy: krótszy i
bardziej czytelny. Zawodowy B2 oznacza, że \emph{rozumiesz} oba i \emph{wybierasz} właściwy.
\end{zasada}

\begin{falle}
Polski styl urzędowy też bywa nominalny, więc kusi, by dosłownie kalkować. Błąd: za długa grupa rzeczownikowa bez czasownika
(\de{die Durchführung der Überprüfung der Bereitstellung \dots}) staje się nieczytelna. Reguła: nie więcej niż dwa dopełniacze
z rzędu; jeśli więcej, podziel zdanie. Drugi błąd: pomylone rodzaje rzeczowników na \de{-ung} (zawsze \de{die}).
\end{falle}

\begin{szwajcaria}
Szwajcarskie umowy, regulaminy i polisy są pisane stylem nominalnym i prawniczym; czytaj je powoli, z słownikiem terminów.
Przykłady: \de{Kündigungsfrist}, \de{Zustellung}, \de{Inkrafttreten}. W pracy ustnej i w e-mailach między kolegami obowiązuje
prostota.
\end{szwajcaria}

\begin{praca}
Komunikat o wdrożeniu: \de{Die Inbetriebnahme des neuen Systems erfolgt am 15.\ Mai. Die Bereitstellung der Schulungsunterlagen
wird vorab sichergestellt.} Ten sam tekst w stylu czasownikowym: \de{Wir nehmen das neue System am 15.\ Mai in Betrieb. Wir stellen vorher die
Schulungsunterlagen bereit.}
\end{praca}

\section{Dialogi}
\textbf{Dialog 1 --- zebranie dotyczące regulaminu}
\begin{dialog}
\turn{A}{Der Entwurf sieht vor, dass die Überprüfung der Anträge innerhalb von fünf Tagen erfolgt.}{Projekt przewiduje, że weryfikacja wniosków następuje w ciągu pięciu dni.}
\turn{B}{Das ist zu lang. Wir sollten eine schnellere Bearbeitung ermöglichen.}{To za długo. Powinniśmy umożliwić szybsze przetwarzanie.}
\turn{A}{Dann treffen wir die Entscheidung heute und stellen die Ressourcen zur Verfügung.}{Zatem decyzję podejmiemy dziś i udostępnimy zasoby.}
\end{dialog}

\section{Ćwiczenia}
\begin{cwiczenia}
\exitem{Nominalizuj: \de{Wir prüfen den Antrag.}}{\de{Die Prüfung des Antrags erfolgt (durch uns).}}
\exitem{Nominalizuj: \de{Man behebt den Fehler.}}{\de{Die Behebung des Fehlers erfolgt.}}
\exitem{Rozpisz: \de{Nach der Einführung des Systems \dots{}}}{\de{Nachdem man das System eingeführt hat, \dots}}
\exitem{Funktionsverbgefüge: \de{entscheiden}.}{\de{eine Entscheidung treffen}}
\exitem{Funktionsverbgefüge: \de{bereitstellen}.}{\de{zur Verfügung stellen}}
\exitem{Funktionsverbgefüge: \de{starten (ein System)}.}{\de{in Betrieb nehmen}}
\exitem{Podaj rodzajnik: \dots{} \de{Durchführung}.}{\de{die}}
\exitem{Przekształć na styl czasownikowy: \de{Die Prüfung der Daten erfolgt durch das Team.}}{\de{Das Team prüft die Daten.}}
\exitem{Przetłumacz formalnie: „Po zakończeniu testów nastąpi wdrożenie.”}{\de{Nach Abschluss der Tests erfolgt die Bereitstellung.}}
\exitem{Przetłumacz: „System zostanie uruchomiony 15 maja.”}{\de{Das System wird am 15.\ Mai in Betrieb genommen.}}
\end{cwiczenia}

\section{Trening słownictwa (generowany)}
\drill{37}

\weekend{Weź akapit z umowy, regulaminu lub komunikatu (np.\ SRF, strona gminy) i przetłumacz go ze stylu nominalnego na
czasownikowy, a potem odwrotnie jeden swój mail. Nagraj streszczenie w trzech zdaniach.}
